\documentclass[UTF8,zihao=-4]{ctexart}
\usepackage[a4paper,margin=2.2cm]{geometry}
\usepackage{xcolor}
\usepackage{hyperref}
\usepackage{enumitem}
\usepackage{fontspec}
\usepackage{amsmath}
\IfFontExistsTF{Microsoft YaHei}{\setCJKmainfont{Microsoft YaHei}}{\setCJKmainfont{SimSun}[BoldFont=SimHei]}
\IfFontExistsTF{Microsoft YaHei UI}{\setCJKsansfont{Microsoft YaHei UI}}{}
\IfFontExistsTF{Consolas}{\setmonofont{Consolas}}{}
\definecolor{linkblue}{RGB}{30,80,145}
\hypersetup{colorlinks=true,linkcolor=linkblue,urlcolor=linkblue,pdfborder={0 0 0}}
\setlength{\parindent}{2em}
\setlength{\parskip}{0.25em}
\linespread{1.16}
\setlist{itemsep=0.2em,topsep=0.35em,leftmargin=2.5em}
\pagestyle{plain}

\title{03\_二阶线性递归关系求解}

\author{}

\date{}

\begin{document}

\maketitle

\section*{03 二阶线性递归关系求解}

\subsection*{题目原文}

递归关系求解：

\[a_{n} = r_{1} a_{n-1} + r_{2} a_{n-2}\]

形似斐波拉契数列。

\subsection*{考查类型}

证明和计算题。

\subsection*{解题思路}

这是二阶线性齐次递推。标准方法是写出特征方程：

\[x^{2} = r_{1} x + r_{2}\]

整理为：

\[x^{2} - r_{1} x - r_{2} = 0\]

求出特征根后，再根据根的情况写通解。

\subsection*{两个不同特征根}

若特征根为 $\alpha$ 和 $\beta$，且 $\alpha \ne \beta$，则通解为：

\[a_{n} = A \alpha^{n} + B \beta^{n}\]

其中 \texttt{A} 和 \texttt{B} 由初始条件决定，例如 $a_0$ 和 $a_1$。

\subsection*{重根}

若特征方程只有一个重根 $\alpha$，则通解为：

\[a_{n} = (A + Bn) \alpha^{n}\]

同样用初始条件求出 \texttt{A} 和 \texttt{B}。

\subsection*{斐波拉契数列例子}

斐波拉契数列满足：

\[F_{n} = F_{n-1} + F_{n-2}\]

对应 $r_1 = 1, r_2 = 1$，特征方程为：

\[x^{2} - x - 1 = 0\]

两个根为：

\[\varphi = (1 + \sqrt{5}) / 2\]
\[\psi = (1 - \sqrt{5}) / 2\]

所以：

\[F_{n} = A \varphi^{n} + B \psi^{n}\]

再由初始条件确定 \texttt{A} 和 \texttt{B}。

\subsection*{渐进复杂度判断}

如果两个根的绝对值不相等，且主导根对应系数不为 0，则 $a_n$ 的增长由绝对值最大的根控制。

例如斐波拉契数列中：

\[|\varphi| > |\psi|\]

因此：

\[F_{n} = \Theta(\varphi^{n})\]

\subsection*{答题模板}

\begin{enumerate}
\item 写特征方程 $x^2 - r_1 x - r_2 = 0$。
\item 求特征根。
\item 按不同根或重根写通解。
\item 代入初始条件求常数。
\item 如题目要求渐进阶，比较特征根绝对值。
\end{enumerate}

\subsection*{例题}

\subsubsection*{例题 1}

题目：求解递推：

\[a_{n} = 3a_{n-1} - 2a_{n-2}\]
\[a_{0} = 1\]
\[a_{1} = 3\]

解答：

特征方程为：

\[x^{2} - 3x + 2 = 0\]

分解得：

\[(x - 1)(x - 2) = 0\]

两个不同特征根为 \texttt{1} 和 \texttt{2}，所以：

\[a_{n} = A\cdot 1^{n} + B\cdot 2^{n} = A + B2^{n}\]

代入初始条件：

\[A + B = 1\]
\[A + 2B = 3\]

解得：

\[A = -1\]
\[B = 2\]

因此：

\[a_{n} = 2^{n+1} - 1\]

\subsubsection*{例题 2}

题目：求解递推：

\[a_{n} = 2a_{n-1} - a_{n-2}\]
\[a_{0} = 1\]
\[a_{1} = 2\]

解答：

特征方程为：

\[x^{2} - 2x + 1 = 0\]

即：

\[(x - 1)^{2} = 0\]

特征根 \texttt{1} 是重根，通解为：

\[a_{n} = (A + Bn)1^{n} = A + Bn\]

代入初始条件：

\[A = 1\]
\[A + B = 2\]

解得：

\[A = 1\]
\[B = 1\]

因此：

\[a_{n} = n + 1\]

\subsection*{常见误区}

不要把所有二阶递推都直接写成斐波拉契形式。只有 $r_1 = 1, r_2 = 1$ 且初始条件对应时，才是标准斐波拉契数列。

不要漏掉重根情形。重根时必须出现 $n \alpha^n$ 项。

\end{document}
